
\begin{refsection}

\chapter{Pantallas apagadas y una difícil convivencia: los y las adolescentes}
\label{cap-26}


\section{Introducción}
\label{cap-26.introduccion}
Para los y las adolescentes, la etapa de aislamiento y suspensión de clases en algunos casos se vivió como una pérdida irrecuperable de ciertas experiencias que no se pueden reemplazar. La centralidad de la escuela como institución y dispositivo de sociabilidad, la importancia del grupo de pares (y su crisis o disgregación frente al aislamiento), los cambios rápidos en el cuerpo y las tensiones familiares frente a la convivencia intensificada serán algunos de los ejes a discutir. En este capítulo se presentan cuatro historias de adolescentes con diferentes características sociales y con distintas configuraciones familiares. Si bien estos cuatro relatos pueden inscribirse dentro de la amplia clasificación de clases medias, presentan diferentes características sociales en relación con sus recursos habitacionales y sus composiciones familiares. También incide la presencia (o no) de familiares a cargo con factores de riesgo.

Los y las adolescentes entrevistados presentan las siguientes características: Sofía, de 14 años, vive en Villa Urquiza y cursa la escuela secundaria. Durante el aislamiento convivió con su madre y la pareja de ella en un departamento, en el marco de una familia ensamblada y con ciertos límites habitaciones. Baltazar, de 16 años, reside en Pilar, también cursa el nivel secundario y, si bien sus padres están separados, vivió el aislamiento en una casa propia, lo cual le brindó mayor espacio y autonomía. Julieta, de 17 años, vive en Wilde con ambos padres en un departamento familiar. Su experiencia estuvo marcada por una convivencia nuclear en un contexto urbano consolidado. Por último, Lucas, de 18 años, comenzó su primer año de universidad en plena pandemia. Vive en el barrio de Villa Crespo con sus padres y un hermano mejor, en una casa propia, también en el marco de una familia de clase media alta.

La elección de estos cuatro relatos nos permite aproximarnos a las distintas formas de habitar el encierro en sectores medios de CABA y del Área Metropolitana de Buenos Aires (AMBA). Abordamos de qué manera, en la fase más estricta de aislamiento, dentro de un sector social relativamente homogéneo en sus condiciones de vida, ciertos atributos vinculados al género y la edad configuraron prácticas y percepciones diferenciales en el espacio doméstico, las rutinas educativas y las dinámicas familiares. Algunos de los interrogantes que aborda son los siguientes: ¿qué lugar tuvo la escolaridad en la vida cotidiana de los y las adolescentes en cuarentena? ¿Cuáles fueron los principales medios de vinculación con sus profesores? ¿Cómo fue la convivencia constante con sus familiares? ¿Qué problemas o barreras se presentaron? ¿Qué estrategias desplegaron los actores?


\section{Juventudes y pandemia: desafíos contemporáneos}
\label{cap-26.juventudes-y-pandemia-desafios-contemporaneos}
La noción de juventud como una etapa diferenciada en el ciclo vital es una construcción relativamente reciente en Occidente, desarrollada entre los siglos XVII y XIX, y ha sido objeto de diversas reflexiones en las ciencias sociales \protect\parencite{4196-CHAVES2009}. Margulis y Urresti (2008) señalaron que las primeras conceptualizaciones abordaban a la juventud en singular, vinculándola a una moratoria social, un período intermedio entre la madurez biológica y la social. Este enfoque se aplicaba principalmente a sectores privilegiados y, aunque representaba un avance al diferenciar la juventud de otras etapas de la vida, no lograba captar la complejidad de las múltiples dimensiones, como género, etnia y clase social, que intervienen en su construcción.

El surgimiento de estudios sobre juventud, a nivel global, se inscribe en los profundos cambios económicos, políticos y culturales de la década del cincuenta, que posicionaron a los jóvenes como un problema social y como nuevos actores políticos y sociales antes de ser objeto de análisis científico \protect\parencite{4196-CHAVES2009}. En Argentina, los estudios sistemáticos sobre juventud comenzaron a consolidarse en los años ochenta, en el marco del retorno de la democracia, y en las décadas posteriores se diversificaron las perspectivas analíticas para dar cuenta de las diversas formas de ser joven. Estas investigaciones pusieron énfasis en los ámbitos del trabajo y la educación, así como en las desigualdades estructurales que condicionan las experiencias juveniles.

La juventud es una categoría analítica histórica y relacional, construida desde los modos en que cada grupo social define y vivencia esta etapa de la vida. No es un estado fijo y universal, ni se la puede delimitar solo con criterios etarios o biológicos, sino que resulta una condición social transitoria, que se configura a partir de las interacciones entre jóvenes y los diversos actores sociales con los que se relacionan, como la familia, los pares y el mercado laboral. Esta perspectiva permite analizar cómo las desigualdades de género, etnia y clase moldean las trayectorias juveniles y las prácticas que desarrollan en su vida cotidiana.

A su vez, autores como Reguillo (2000) y Margulis (2001) insisten en que la juventud debe analizarse en el marco del entramado social e histórico que le otorga significación. Esto implica considerar tanto las representaciones sociales, estereotipos y prescripciones sobre lo que significa ser joven, como las prácticas efectivas y las relaciones que los jóvenes construyen. La juventud supone un proceso de individualización en el que los jóvenes comienzan a distanciarse de la familia como un «otro significativo» (Singly, 2000), dando lugar a la construcción de nuevas redes y relaciones relevantes en la formación de su identidad.

Ahora bien, ¿qué pasó con las juventudes durante la cuarentena? ¿Cómo se reconfiguró este proceso de individualización en el marco de una pronunciada convivencia familiar? En este marco, las juventudes experimentaron una disrupción significativa en sus procesos de socialización, construcción identitaria y desarrollo emocional. Como lo destacan UNICEF (2020a) y la Sociedad Argentina de Pediatría (SAP, 2020), el aislamiento social derivó en un aumento del tiempo frente a las pantallas, trastornos del sueño y emociones negativas como angustia, tristeza y miedo. Failla \protect\parencite*{4198-FAILLA2021} subraya que estas condiciones no solo amplificaron el impacto psicológico del aislamiento, sino que también reforzaron el sentimiento de desconexión y desamparo entre los jóvenes. En este contexto, la relevancia de los espacios privados, como la habitación propia, y las redes de amigos como escenarios para desplegar autonomía se hizo particularmente evidente. Sin embargo, para aquellos adolescentes que carecían de un entorno doméstico favorable, el confinamiento exacerbó las tensiones en la familia, limitando las posibilidades de autonomía.


\section{Cada relato, un universo: los perfiles de nuestros entrevistados}
\label{cap-26.cada-relato-un-universo-los-perfiles-de-nuestros-entrevistados}
Julieta tenía 17 años en 2020 y vivía con sus padres en un departamento de Wilde. Le tocó transitar el final de la secundaria durante la cuarentena: el cansancio y la desconexión con las clases virtuales la llevaron a abandonar las plataformas digitales. Luego comenzó a prepararse para ingresar a la carrera de Medicina y sintió la falta de preparación, que vincula con los contenidos que no vio durante la pandemia. Durante ese tiempo diferente, desarrolló nuevas aficiones como el maquillaje y los videojuegos, y vivió tensiones con su grupo de amigos, llegando incluso a pelearse con alguno de ellos. Las salidas clandestinas, las videollamadas diarias con su familia y el sentimiento de culpa la acompañaron durante un largo período, además de la tristeza por la pérdida de su abuela.

Lucas tenía 18 años y vivía con sus padres –ambos profesionales– y un hermano menos en una casa en el barrio de Villa Crespo. Cursó cuarto y quinto año del secundario con buena parte de las clases a distancia y recuerda esos días como una etapa de angustia y frustración. La falta de contacto con sus compañeros lo afectó mucho, a pesar de los esfuerzos por mantenerse conectado a través de redes sociales. Las tensiones en casa, debido a que su padre era de riesgo, le produjeron ansiedad y problemas para dormir. Actualmente acude a terapia para procesar lo vivido, consciente de que las relaciones personales y la salud mental no volvieron a ser las mismas tras el confinamiento.

Baltazar tenía 16 años en 2020, vivía con su madre en Pilar, en una casa propia. Sus padres estaban separados. Atravesó la pandemia entre videojuegos, series y ajedrez. Aunque logró conectarse mejor con sus compañeros durante las clases virtuales, no deja de preguntarse cómo habría sido su experiencia de haber tenido un contacto presencial. Siente que el encierro le quitó meses de encuentros con amigos, pero por otro lado mejoró su relación con la tecnología y el entretenimiento online.

Por último, Sofía tenía 12 años en 2020 y convivía con su mamá. Como sus padres estaban separados, cada tanto veía a su padre. Le tocó comenzar la secundaria en plena pandemia y se enfrentó al desafío de conocer a un grupo de pares nuevo a través de las pantallas de Zoom. La vergüenza de encender la cámara y participar la paralizaba, lo que la llevó a refugiarse en nuevas aficiones como el dibujo y la música. Su mamá seguía trabajando a la distancia, padecía un cuadro de ansiedad que se agravó con la pandemia y esto a su vez afectó a Sofía.

Mediante las historias y los testimonios de los adolescentes analizaremos los cambios en las rutinas escolares durante la pandemia, las dinámicas familiares y las estrategias que cada uno de ellos desplegó durante el aislamiento.


\section{Nuevas configuraciones escolares: la escuela dentro de la casa}
\label{cap-26.nuevas-configuraciones-escolares-la-escuela-dentro-de-la-casa}
En la pandemia, la vida cotidiana de los adolescentes se vio muy trastornada. Las fronteras entre lo público y lo privado dentro del hogar se desdibujaron. La vivienda se convirtió en escenario de múltiples actividades: tareas domésticas, trabajo remunerado, educación, ocio, vínculos afectivos y nuevas experiencias ligadas a la salud física y mental. Aparecieron preocupaciones por el contagio propio o de los seres queridos, y también formas de miedo, encierro y sobrecarga emocional. Aunque todos los sectores sociales se vieron afectados por el aislamiento, no todos lo vivieron de la misma manera.

En el caso de los y las adolescentes, que no fueron el grupo más expuesto al riesgo sanitario, los efectos aparecieron en otros planos. Según Cardinaux \protect\parencite*{4191-CARDINAUX2020}, este sector deberá enfrentar, en el corto y mediano plazo, consecuencias duraderas en términos de oportunidades truncadas. A esto se suma una suposición extendida durante la pandemia: al considerarlos más autónomos que los niños pequeños, se creyó que podían gestionar por sí mismos sus aprendizajes. Sin embargo, ese supuesto no se cumplió.

La suspensión de la presencialidad alteró las reglas que organizaban la vida escolar. La rutina institucional, que regulaba la asistencia, la promoción y la certificación, perdió fuerza. Ante esa disrupción, los adolescentes tuvieron que reorganizar su tiempo, sus vínculos y sus formas de aprender, muchas veces sin acompañamiento suficiente.

La pandemia cambió la «domiciliación» de lo escolar, trasladándola hacia el espacio doméstico \protect\parencite{4197-DUSSEL2020}. Este cambio creó tensiones difíciles de resolver y a su vez reveló la importancia de la estructura material y simbólica de la escuela para producir un \emph{espacio otro} donde se despliega la autonomía intelectual y afectiva progresiva. ¿Qué pasa con la mudanza de la escuela al espacio doméstico? ¿Qué sucede con la pérdida de la diferenciación de roles, reglas propias de la escuela y su reacomodación al ámbito hogareño? ¿Qué nuevas reglas y autoridades se instalan?


\begin{gbcita}
«Yo sufrí bastante la virtualidad. Había algo que me era muy tosco y lento de tener clases de esa manera. Yo sé que aprendí muy poco en el terreno académico estos dos últimos años, o sea, no digo que no aprendí nada, pero fue como una regresión muy fuerte por todo el tema de que estás encerrado en tu casa, viendo una pantalla por varias horas al día, y la cabeza te queda agotada» (Lucas, 18 años).

«Era feo cursar virtual. No sé si en mi caso fue tanto por no entender algo, sino porque era raro usar zoom. Porque si vos ves presencialmente al profesor, le preguntás cosas, pero por zoom no se lo preguntás porque todo queda grabado, y al menos a mí me daba cosa» (Sofía, 14 años\textbf{).}

«Me despertaba a la mañana y tenía que estar en el zoom muchas horas por día. Y podía estar desde mi cama escuchando la clase» (Julieta, 17 años).

«Yo vi gente que tenía clases todos los días como un horario normal de colegio, pero yo no tuve tantas. Por ejemplo, tenía una profesora de inglés que mandaba actividades y no hacía clases virtuales. Nos mandaban la tarea por una plataforma y yo hacía todo el lunes y después tenía toda la semana libre» (Baltazar, 16 años).

\end{gbcita}

Más allá de las edades y el género de los entrevistados, en todos los relatos podemos observar que hay una clara ruptura de los espacios, las rutinas diarias tal y como las conocían, y una incomodidad por tener que cursar mediante una pantalla. Algo que aparece en las cuatro entrevistas es que sus experiencias educativas fueron clases sincrónicas mediante la plataforma zoom. Algunos de los entrevistados tuvieron pocas, como el caso de Baltazar, que nos cuenta que tuve menos de lo que él creía en relación con la cantidad de materias que cursaba. Pero los cuatro relatos transitaron su escolaridad mediante la sincronicidad.\footnote{Un informe de la Universidad Torcuato Di Tella (2020) argumenta que
      las diferencias de acceso a la tecnología y conectividad se relacionan
      con el despliegue de dos modelos de escolarización bien diferenciados:
      escuelas WhatsApp y escuelas Zoom. El documento expone que, en los
      contextos más desfavorecidos, los docentes utilizaron la aplicación de
      WhatsApp como principal herramienta, y en segundo lugar el correo
      electrónico: dos instrumentos de naturaleza asincrónica y fácil
      acceso. En cambio, en los contextos con población favorecida se
      utilizó como principal herramienta las videoconferencias sincrónicas
      (la más extendida es Zoom), con una frecuencia de intercambio diario
      en la mayoría de las escuelas. En ese sentido, las cuatro experiencias
      se enmarcan en un sector social relativamente homogéneo en sus
      condiciones de vida.} Es decir, la virtualidad fue una estrategia clave para suplir la falta de escolaridad presencial. Sin embargo, los adolescentes entrevistados manifestaron sensaciones de malestar en esas clases virtuales debido a la baja interacción con docentes y pares, así como sentimientos de sobreexposición y vergüenza (Moguillansky \emph{et al}, 2022).

Al respecto, Inés Dussel \protect\parencite*{4197-DUSSEL2020} sostiene que la educación a distancia generó cierta individualización del trabajo pedagógico. La escuela es un espacio colectivo, un espacio de lo común; y la falta del aula presencial produce una ausencia de otras voces que hace que se pierda algo importante. En el relato de Sofía podemos vislumbrar su saturación ante tanta observación y atención adulta, a tal punto que le da vergüenza hablar porque «todo queda grabado». Este tipo de experiencias pone en evidencia que en el aula los estudiantes se reparten la carga, se escuchan, aprenden de lo que dicen otras u otros, encuentran con quién o en dónde esconderse de la demanda adulta. En ese sentido, Dussel subraya que la nueva escena pedagógica virtualizada tiene algo de panóptico, y que, lejos de convertirla en un asunto público, la convierte en un juego de exigencias y de espejos que desplazan que se está ahí para enseñar y aprender y no para atender a demandas individuales \protect\parencite[p. 5]{4197-DUSSEL2020}.


\begin{gbcita}
«La mayoría de los docentes no daban clases todo el tiempo. Algunos dejaban un mensaje en el \emph{classroom}. Y eso a veces daba la sensación de que no existían. Pero cuando volvimos a las clases se notaba el esfuerzo que intentaban hacer para enseñarnos» (Baltazar, 16 años).

«Al principio de 5to, cuando recién empezó la cuarentena, tenía ganas de hacer las tareas. Todo el primer trimestre me fue re bien porque hacía todo. Me despertaba a la mañana, me conectaba a los zooms, y hacía las tareas. Tenía una rutina como si estuviera dentro de la escuela. Pero ya para el segundo trimestre no hice nada. Y en sexto me pasó lo mismo. Entendí que no importaba si hacía todo o no hacía nada, igual iba a aprobar. Y entonces se me fueron las ganas de intentar» (Julieta, 17 años).

«Muchas clases me resultaban muy aburridas. Era muy pesado levantarte a la mañana y escuchar a una profesora hablar por dos horas seguidas de lo mismo» (Sofía, 14 años).

\end{gbcita}

Tal como podemos observar, la tarea escolar que recibieron los y las adolescentes a través de tecnologías digitales tomó distintas connotaciones. En algunos casos, resultaba una organizadora estructural del día, pero para otros aparecía como excesiva, o como una actividad rutinaria y obligatoria, que limitaba la función de la escuela a enviar la tarea y la de los estudiantes a realizarla. Julieta comparte una experiencia de desmotivación progresiva. Al inicio de la cuarentena, se sentía comprometida, con una rutina que le daba sentido a las tareas, pero con el tiempo, su interés disminuyó al darse cuenta de que el esfuerzo no influía en sus calificaciones. La monotonía de las clases virtuales y la falta de motivación en los contenidos también generaron desánimo. La mayoría de los testimonios evidencian la dificultad de mantener el interés y la conexión emocional en un formato educativo distante y repetitivo. Pero el límite no es solamente el de la capacidad de las plataformas y las emociones que los y las adolescentes tienen al respecto, sino también el de la atención, ya que resultaba arduo tener que estar varias horas conectados.

Este tipo de experiencias no hacen más que visibilizar los problemas de la nueva economía de la atención. Al respecto, Dussel sostiene que en el último tiempo la atención es pensada como mercancía. Plataformas mundiales intentan captar segundos o minutos de nuestra atención en avisos publicitarios. Si bien esto no es una novedad, el impacto y el avance es cada vez mayor. En esta nueva economía de la atención, es difícil concentrarse en algo. Y si en la escuela concentrarse en el estudio ya era difícil, en el ámbito doméstico lo fue mucho más porque cobran otra dimensión las interrupciones continuas de lo doméstico en las clases o la posibilidad de poder apagar una cámara y seguir estando al mismo tiempo:


\begin{gbcita}
«Yo me conectaba a las clases, pero sin la cámara. Nadie la tenía prendida. Yo no la prendía porque me daba vergüenza. Igual pensaba, “pobre profe está viendo la pantalla negra”, porque era todo silencio y oscuridad, y cuando hablaba estaban todos callados. Porque si te confundías quedabas como muy expuesto. Y nadie contestaba. La profesora tenía que decir el nombre específico de la persona. Y había veces que ni así contestaban… decían que se les había roto el micrófono, y que por alguna razón tampoco podían escribir» (Sofía, 14 años).

«Yo prendía la cámara, era de los pocos. Porque me daba un poco de cosa; es que vos veías a la mayoría de la gente y nadie con la cámara prendida, era la nada misma y nadie respondía y la mayoría estaba desde la cama» (Lucas, 18 años).

\end{gbcita}

La virtualización de las clases durante la pandemia trajo consigo nuevas experiencias escolares que se cruzaron con diferencias de género. Una de ellas puede observarse en el uso de la cámara durante las clases sincrónicas. En el testimonio de Sofía aparece con claridad una sensación de vergüenza frente a la posibilidad de encenderla y participar. Esa incomodidad se relaciona con formas de exposición que, en la adolescencia, adquieren un peso particular. Los cambios corporales, junto con la presión social sobre la apariencia, suelen afectar con más intensidad a las mujeres. Esta situación invita a una pregunta más amplia: ¿por qué la cámara encendida durante la clase genera más vergüenza que la exposición en redes sociales, donde las imágenes circulan sin mediación adulta y con una posibilidad real de recibir críticas? La presencia del docente y del grupo escolar, en un contexto que mezcla lo público y lo íntimo, podría amplificar la sensación de desajuste, sobre todo cuando el aula virtual se superpone con el espacio del hogar.

El encierro también volvió más visible ese otro plano de exposición: el fondo de la casa, la habitación compartida, los sonidos familiares. Para muchos adolescentes, mostrar su entorno doméstico resultó incómodo. El temor a ser juzgados por el espacio donde viven reforzó esa incomodidad. Giddens \protect\parencite*{4201-GIDDENS1997} sugiere que las situaciones de crisis como la pandemia pueden quebrar la seguridad ontológica, es decir, la confianza básica que permite actuar en el mundo. En ese marco, emergen sensaciones de inadecuación y vergüenza, que debilitan la autonomía y afectan el modo en que los sujetos se presentan ante los demás.


\section{Rutinas alteradas: cambios en el sueño y el impacto del sedentarismo en los y las adolescentes}
\label{cap-26.rutinas-alteradas-cambios-en-el-sueno-y-el-impacto-del-sedentarismo-en-los-y-las-adolescentes}
Durante la pandemia, la vida cotidiana de los adolescentes se volvió más sedentaria. En las zonas urbanas de la Argentina, al igual que en muchos otros países, se observó una disminución en la actividad física y un uso más prolongado de pantallas (Páez y Simoni, 2022). La interrupción de las clases presenciales, las restricciones para circular y el cierre de espacios recreativos redujeron las oportunidades para moverse, jugar o entrenar. Según la Sociedad Argentina de Pediatría, los adolescentes deberían realizar al menos una hora diaria de actividad física de intensidad moderada. Sin embargo, durante el confinamiento, pocos lograron sostener ese mínimo. La falta de movimiento no solo afectó la salud, sino también el estado de ánimo, el sueño y la capacidad de concentración, aspectos claves en una etapa de crecimiento y transformación como la adolescencia. En la pandemia esto disminuyó enormemente:


\begin{gbcita}
«Lo de mi hermano fue un desastre porque le pegó muy rara la cuarentena. En ese momento tenía 14 años. Estaba todo el día encerrado con la computadora. Yo lo iba a visitar al cuarto cada tanto a ver cómo andaba porque no lo veía nunca. A veces aparecía para almorzar, pero comía re rápido y se iba corriendo al cuarto. Estaba encerradísimo y era una madriguera el lugar. Sin luz, toda ropa tirada, todo hecho un desastre, comida en el escritorio, y él mirando la computadora. Si yo le decía “che, ¿no me acompañas a comprar una lata de coca al quiosco?”, y no, no salía. Siempre le gustó jugar los jueguitos en la compu, pero nunca fue este nivel de encierro. Tenía todo el día las persianas cerradas, no veía nadie, hablaba con los amigos por celular y no salía del cuarto» (Lucas, 18 años).

«Me pasó de estar todo el día encerrada, con la ropa con la que me despertaba. A veces no me bañaba o no me cambiaba de ropa en todo el día» (Sofía, 14 años).

\end{gbcita}

Estos fragmentos indican el profundo impacto que el encierro tuvo en las rutinas diarias de los y las adolescentes, en particular en lo relacionado con el sueño y el sedentarismo. La cita de Lucas describe cómo su hermano pasó de una relación más equilibrada del uso de la computadora a un aislamiento extremo, donde la actividad física y los vínculos con su entorno fueron drásticamente reducidos. Este comportamiento no solo ilustra la falta de actividad física recomendada por organismos de salud, sino también un deterioro del espacio personal, que se convierte en una «madriguera» desordenada y oscura, símbolo del aislamiento emocional y social. Por otra parte, Sofía describe una situación similar en la que su rutina básica, como cambiarse de ropa o salir de la cama, también quedó afectada. Esto subraya cómo el encierro forzado afectó las corporalidades adolescentes, que se vieron inmovilizadas y limitadas a comportamientos sedentarios. La falta de interacción física en espacios sociales como la escuela, clubes, bares, sumado al exceso de tiempo frente a pantallas, alteró la salud física y psicológica, dejando a muchos jóvenes a la «intemperie», desprovistos de los cuidados y contención habituales (Páez y Simoni, 2022). A la vez, este sedentarismo trajo cambios corporales. A la salida de la pandemia, con el regreso a clases presenciales, el reencuentro entre los adolescentes fue complejo.


\begin{gbcita}
«Muchos de mis compañeros estaban muy cambiados. Hay un compañero que estaba así (gesto con las manos) de alto, como dos metros (risas). A mí me daba cosa verlos y que ellos me vean a mí porque no me cuidé tanto físicamente. A todos nos pasó lo mismo, me acuerdo de que nos daba vergüenza sacarnos el barbijo en las clases de educación física. Era raro ver la cara entera de alguien. Ahora aprovecho el barbijo para esconderme un poco también» (Sofía, 14 años).

«La vuelta a clase en 5to año fue darme cuenta de que realmente no reconocía a un montón de gente. Personas que tuve un vínculo cercano por tres años y de la nada tuvieron todo un cambio en el que no pude estar presente, nadie pudo estar presente, y pasaron cosas. Darte cuenta de que toda la gente que conocías con 15 o 16 años, ahora tienen 18 y son personas completamente distintas, que les pasaron mil cosas, y no te vieron por un año, eso también fue muy fuerte» (Lucas, 18 años).

\end{gbcita}

Sofía menciona que el reencuentro con sus compañeros luego de la pandemia fue marcado por cambios físicos notorios, tanto en ella como en los demás. La referencia a un compañero que creció significativamente durante ese período expone lo rápido que sus cuerpos cambiaron en la ausencia del contacto cotidiano. Por otro lado, el uso del barbijo no solo se transformó en una medida sanitaria, sino también en una herramienta para ocultar su apariencia e inseguridades que surgieron durante el aislamiento. En lo referido a los cambios, Lucas experimentó una fuerte desconexión con sus compañeros con quienes compartió vínculos cercanos antes de la pandemia. Su relato sostiene que se convirtieron en «personas completamente distintas» luego de esta experiencia. En estos dos fragmentos podemos observar que el aislamiento extendido no solo afectó sus cuerpos, sino que también alteró las dinámicas sociales y afectivas, donde pasaron cosas sin que nadie estuviera presente para compartirlas.

En cuanto a las rutinas de sueño, los y las adolescentes necesitan dormir más que los adultos: al menos ocho por día. El descanso es vital para desempeñarse en la escuela y en la vida. Los adolescentes privados de sueño no pueden concentrarse y su aprendizaje y sus calificaciones se ven afectados. También están más irritables y ansiosos a lo largo del día, por lo que sus relaciones se resienten \protect\parencite{4200-HAIDT2024}. En el caso de Baltazar, recuerda que durante la pandemia se trastornaron sus horarios y tuvo dificultades para dormir:


\begin{gbcita}
«Mis horarios estaban todos cambiados. Yo me acuerdo de que me levantaba muy tarde, había días que me iba a acostar a las 7 de la mañana y me levantaba a las 3 de la tarde. Después me costó un montonazo volver a levantarme temprano y empezar a cumplir el horario normal. A veces decía “bueno me voy a ir a dormir temprano”, y no podía. Siempre me quedaba despierto porque vivía en llamada con mis compañeros y no te dabas cuenta cómo pasaba el tiempo y miraba el reloj y eran las 6 de la mañana» (Baltazar, 16 años).

\end{gbcita}

En contraste, Sofía cuenta que durante el aislamiento mantuvo una rutina con horarios más o menos similares a los previos, ordenando el día en función de las clases virtuales, aunque algunos días se levantaba más tarde:


\begin{gbcita}
«Mantenía la misma rutina porque tenía las clases virtuales. Por ahí los días que no tenía clases virtuales me despertaba más tarde» (Sofía, 14 años).

\end{gbcita}

El testimonio de Baltazar nos muestra cómo los horarios irregulares y la exposición a pantallas hasta altas horas de la madrugada afectaron su patrón de sueño. La descripción de sus noches largas, hablando con sus amigos por teléfono, sumadas a la dificultad para reestablecer una rutina normal después del confinamiento, nos indica cómo el cambio en las rutinas diarias impactó negativamente en su capacidad para mantener horarios regulares y en su bienestar general. Su experiencia pone de relieve cómo las distracciones digitales contribuyen a la alteración del sueño que afecta su capacidad para concentrarse y participar en actividades diarias. Por otro lado, el relato de Sofía revela una adaptación parcial a la nueva rutina con clases virtuales, que permitió cierto grado de flexibilidad en sus horarios de sueño. El comentario sobre los cambios en el sueño vinculado a la ausencia de escolaridad presencial destaca cómo la falta de una rutina estructurada y la flexibilización de horarios contribuyeron a una mayor relajación en las rutinas de sueño.

En suma, la interrupción de las rutinas escolares y la falta de interacción social afectaron el bienestar emocional, pero también su relación con el cuerpo, la alimentación y los cuidados en casa. La combinación de estos factores trastocó sus vínculos con los demás y con ellos mismos y exacerbó el malestar y la sensación de desamparo que muchos experimentaron en este período (Tuñón, 2022).


\section{Cambios en las prácticas de cuidado durante la pandemia}
\label{cap-26.cambios-en-las-practicas-de-cuidado-durante-la-pandemia}
Tradicionalmente, el cuidado fue concebido como una actividad unidireccional, donde los adultos y las instituciones eran responsables de proteger a los más jóvenes, considerados vulnerables. Sin embargo, la pandemia del COVID-19 introdujo un cambio significativo en esta noción. Según Barcala (2020), los adolescentes no solo se percibieron como receptores de cuidados, sino también como cuidadores activos. Este cambio se evidenció en los y las adolescentes cuando comenzaron a asumir roles de cuidado hacia sus familiares y la sociedad en general:


\begin{gbcita}
«Mi mamá siempre tuvo traumas de la infancia. Durante la pandemia tenía mucha ansiedad y estrés. No era violenta ni nada. Pero ella estaba mal, y verla así me ponía mal. Ella estaba todo el día en la computadora y no hablaba con nadie. A veces estaba en el cuarto de ella como para acompañarla, pero no podíamos tener una conversación porque tenía que estar constantemente concentrada en la computadora» (Sofía, 14 años).

«Mi vieja labura de arquitecta, así que en ese momento laburo no tenía, y mi viejo es emprendedor y labura mucho con acciones, pasa mucho tiempo en casa. Los dos tenían poco trabajo y era muy intensa la convivencia 24/7. Hubo mucha pelea, la verdad que fue complicado. Cuando estás conviviendo con alguien tanto tiempo y tan pegado, es imposible no tener roces. Encima mi hermano y yo somos adolescentes, claramente nos vamos a pelear» (Lucas, 18 años).

\end{gbcita}

En ambos relatos se observa la experiencia de roles de cuidado activos. Sofía, por ejemplo, describe cómo la salud mental de su madre le afectaba y cómo ella intentaba acompañarla, incluso cuando no podían tener una conversación por el trabajo. En este caso, Sofía no solo experimentó el impacto emocional de ver a su madre en una situación de estrés, sino que también adoptó una postura de apoyo, aunque sea de manera silenciosa. Por otro lado, Lucas describe una dinámica familiar más tensa, marcada por el confinamiento. La cercanía obligada por la pandemia exacerbó los conflictos familiares. La repetición de roces familiares, las peleas constantes y el resentimiento acumulado indican cómo el confinamiento generó situaciones de agotamiento emocional en los y las adolescentes:


\begin{gbcita}
«Al principio las compras fueron un desastre. Me acuerdo de que a mi viejo le agarró la paranoia. Era de esos que veían mucho las noticias, se perseguía y todo era terrible. Entonces las compras eran una operación complejísima y todos limpiábamos cada cosa que comprábamos con alcohol y con lavandina. Fue muy difícil la convivencia con mi familia. Tener que vivir con una persona que es población de riesgo y que encima es de base muy hipocondríaca, no es para cualquiera. La situación en casa fue muy complicada y hubo muchas peleas y mucho resentimiento» (Lucas, 18 años).

\end{gbcita}

El miedo al otro se convirtió, durante la pandemia, en un factor que favoreció el encierro y debilitó los lazos con el entorno. En muchas familias, este temor provocó tensiones difíciles de manejar. Según Segato (2018), este tipo de reacciones no solo reflejan una preocupación por el bienestar familiar, sino que también muestran cómo el aislamiento afectó el tejido social. El testimonio de Lucas sobre las compras ilustra con claridad cómo ese temor al contagio se transformó, en su hogar, en un comportamiento obsesivo por parte de su padre, intensificando la tensión familiar. Ese miedo desbordado a la infección puede interpretarse como una forma extrema de cuidado hacia los propios, pero también como un factor que profundizó el aislamiento y debilitó los lazos con el entorno social. A pesar de los discursos que señalaron a los adolescentes como indiferentes frente a la crisis, la experiencia mostró otra cosa. En distintos hogares, los jóvenes se cuidaron entre sí y cuidaron a los suyos. Muchos expresaron angustia por la posibilidad de contagiar a personas queridas y sintieron el peso de una responsabilidad que no eligieron (Moguillansky \emph{et al}., 2022).

En algunas investigaciones se advierte un cambio en la distribución de las tareas del hogar. Ianina Tuñón y su equipo (2021) observaron que, al reducirse el tiempo dedicado a la escuela, los adolescentes asumieron nuevas responsabilidades domésticas. Lejos de la imagen de pasividad, este involucramiento en el cuidado cotidiano muestra una forma concreta de compromiso con la vida familiar. En muchos casos, ese esfuerzo se acompañó de sentimientos de culpa y sobrecarga, sobre todo cuando había personas del grupo de riesgo en la casa.


\section{El hogar como refugio: desafíos de volver al mundo exterior}
\label{cap-26.el-hogar-como-refugio-desafios-de-volver-al-mundo-exterior}
Entendida como epicentro de la vida privada, la casa es la unidad productiva y reproductiva en torno de la cual se organizan buena parte de las tareas y actividades que conforman la vida cotidiana. Más allá de las innegables cualidades simbólicas que unen al habitante con su espacio vivido, la casa está asociada con la noción de «morada» desde la cual, entre otros procesos, se garantiza la reproducción de la fuerza de trabajo a través del alimento, el descanso, la higiene, etcétera \protect\parencite{4199-GARDINER1997}. Al mismo tiempo, constituye un lugar que hace posible el despliegue y resguardo de la intimidad, los afectos y proyectos. Como tal, la casa ofrece abrigo y sostén, presentándose a sus habitantes como garantía del cuidado y protección necesarios para la vida. Al sintetizar las conexiones más estrechas entre espacio e intimidad, la casa es un centro espacial y afectivo a partir del cual se producen, fijan y evalúan los contenidos de las relaciones de pertenencia, identidad, extranjería, etcétera, que organizarán las más diversas interacciones sociales. En el mejor de los casos, bajo el cobijo de la casa los sujetos inician sus aprendizajes sociales y afectivos; aprendizajes que irán forjando el andamiaje de costumbres, hábitos y valores que marcarán sus trayectos por el mundo social (Hochschild, 2008; Pautassi y Zibecchi, 2013). Más allá de que la casa cobró una importancia relevante durante el confinamiento, los y las adolescentes estuvieron de acuerdo en que su lugar «preferido» dentro de ese entorno fue su habitación:


\begin{gbcita}
«Bueno mi casa es un terreno en el que hay tres casas. Y en el mismo terreno viven mi tía y mi abuela. Y muchos animales que son de ellas, gatos y perros. El espacio que más usé fue mi pieza, básicamente vivía acá. Me acuerdo de que a veces pasaban tres días en que no salía, me la pasaba jugando a los jueguitos en mi cuarto» (Baltazar, 16 años).

«Tengo 14 años, vivo en la casa de mi papá y en la casa de mi mamá. Ahora estoy la semana con mi mamá porque el colegio está cerca de la casa de mi mamá, y los fines de semana con mi papá. En el espacio que más estoy es en mi cuarto» (Sofía, 14 años).

«Después de la pandemia, quedarme en casa es algo que hago mucho. Puedo estar 5 días seguidos en mi casa y no me parece raro ni nada, cuando en otro momento quizás sí» (Julieta, 17 años).

\end{gbcita}

Comprendida en forma global, la habitación es un «pequeño universo» en el que se articulan y transcurren distintas actividades personales. En todos los casos, los adolescentes destacan que su habitación es el espacio que más utilizan y en el que se sienten más cómodos. Esta elección indica que la habitación no solo representa un lugar físico, sino también un espacio personal y emocional significativo para ellos. Baltazar, por ejemplo, vive en un terreno compartido con familiares y animales, pero su habitación es su lugar predilecto, lo que indica una necesidad de un espacio propio en un entorno con muchas personas. Por otra parte, Sofía menciona que su habitación en la casa de su mamá es su espacio preferido mientras está allí. Esto indica que, a pesar de la dualidad de sus entornos, su cuarto le proporciona una sensación de estabilidad y pertenencia. Julieta, en cambio, experimentó un cambio luego de la pandemia ya que ahora le resulta normal pasar varios días seguidos en su casa. Este relato pone de manifiesto las transformaciones en las dinámicas de vida cotidiana que trajo el aislamiento, en particular en relación con la percepción del tiempo que pasamos en el hogar.

Según Segura y Caggiano \protect\parencite*[pp. 6-7]{4190-CAGGIANO2021}, al alterarse los tiempos de habitación de los espacios de la casa, el trabajo, la distensión y el esparcimiento convergen en un mismo lugar, lo que en la experiencia de los y las adolescentes permitió dar continuidad de otra manera a las actividades que dejaron de realizarse fuera del hogar. En efecto, «el cuarto propio», donde mayormente estuvieron recluidos, fue el espacio doméstico central en los relatos de los adolescentes entrevistados, inclusive entre aquellos que manifestaban tener una casa grande o ambientes al aire libre como un patio o balcón.

Ahora bien, si en la pandemia la casa era entendida como un refugio, ¿qué pasa cuando se vuelve «a la normalidad»? En las entrevistas se registra una sensación de tiempo perdido y una nostalgia por los vínculos y experiencias que parecieran no poderse recuperar:


\begin{gbcita}
«Yo siempre me pregunto cómo hubiera sido el 2020 si íbamos al colegio. Seguro que pasaba más tiempo con mis amigos. Creo que eso fue una de las peores cosas, el no tener interacciones con los amigos en persona porque literalmente no los vi por un año» (Lucas, 18 años).

«Ir al colegio para mí era vivir la experiencia de estar en quinto. Más que nada la ilusión de que el año que viene iba a ser promo y que quinto es como el año más tranquilo, porque no tenés que elegir carrera ni nada de eso y podés disfrutarlo. Y bueno, no lo pude vivir» (Baltazar, 16 años).

«Egresé el año pasado del colegio y ahora estoy haciendo el curso de ingreso para entrar a la facultad de Medicina. Es un mes de curso y te toman una prueba final. Estoy bastante nerviosa porque es muy difícil. Aparte no conozco a nadie. Dentro de lo que me preocupa el curso de ingreso, está también eso de si voy a ser amigos o no, si voy a poder socializar» (Julieta, 17 años).

\end{gbcita}

Una coincidencia entre los adolescentes es la sensación de que perdieron un tiempo y/o una experiencia específica e irrecuperable durante la pandemia. Este sentimiento acrecienta emociones como la angustia, el miedo o la tristeza. La interrupción de rituales de pasaje, como los últimos años de escuela o la transición a la universidad, revelan la importancia de estos momentos para la formación de la identidad y la sociabilidad. El aislamiento y la falta de interacciones presenciales generaron ansiedad, como el caso de Julieta, sobre la capacidad para socializar y establecer nuevas conexiones en un mundo postpandémico.

Por último, el regreso a la escuela y el reencuentro con sus pares está asociado a sensaciones de vergüenza y desconocimiento. Los demás cambiaron y muchos indican la dificultad para recuperar viejas amistades o, directamente, perderlas del todo y entrar en la trabajosa construcción de nuevos lazos. Esta sensación aparece con más fuerza en aquellos jóvenes que estaban en las etapas iniciales y finales de su educación secundaria al momento del aislamiento más estricto (Moguillansky \emph{et al}, 2022).


\section{Conclusiones}
\label{cap-26.conclusiones}
La pandemia marcó a una generación de estudiantes que atravesaron años cruciales de su vida académica y social en aislamiento, donde cada uno enfrentó (y sigue enfrentando) desafíos y realidades distintas. La suspensión de las clases presenciales y la necesidad de reorganizar las tareas de cuidado transformaron el hogar en un espacio multifuncional, donde los distintos miembros de la familia debieron gestionar sus actividades cotidianas a través de dispositivos conectados a internet \protect\parencite{4189-ACEVEDO2021,4192-ANDERETESCHWAL2020,4182-HEREDIA2022}.

Sin embargo, la experiencia de la pandemia no fue vivida de manera homogénea. El análisis de las entrevistas muestra similitudes y diferencias en las vivencias de los adolescentes durante el confinamiento. En cuanto a los puntos en común, todos experimentaron rupturas en su rutina académica y social, lo que generó sentimientos de aislamiento y la necesidad de buscar nuevas aficiones. Sin embargo, las diferencias radican en la manera particular en que cada uno procesó este contexto adverso: mientras Julieta y Lucas experimentaron un impacto emocional profundo, marcado por el malestar escolar y el acompañamiento terapéutico, Baltazar encontró cierto equilibrio en el entretenimiento digital, y Sofía se refugió en actividades creativas. Lo interesante de estas discrepancias es que los adolescentes no se limitaron a repetir los discursos de los adultos para entender estos cambios. Desarrollaron sus propias elaboraciones personales con los recursos de cada uno.

Por último, un aspecto relevante que merece un análisis más profundo se relaciona con el retorno a las prácticas de sociabilidad presenciales y a los espacios de interacción cara a cara, dado que los contextos a los que se regresa parecen haber sufrido transformaciones significativas. También sería interesante analizar si en este regreso hubo diferencias de experiencias según el género, ya que en el análisis de las entrevistas vislumbramos como Sofía, por ejemplo, narró que tenía vergüenza debido a la exposición frente a pantallas.

En conclusión, la pandemia no solo modificó las dinámicas familiares, sociales y educativas, sino que también reconfiguró las formas en que los adolescentes se relacionan con el mundo y consigo mismos. ¿Qué nuevas formas de habitar el espacio público y privado surgirán a partir de estos cambios? ¿Cómo se verán afectadas las relaciones intergeneracionales, los vínculos sociales y la percepción de los jóvenes sobre su futuro? Sin lugar a duda, son interrogantes que aún no pueden ser respondidos, producto del contexto emergente y cambiante que transitamos. Pero plasmarlos, es un buen punto de partida.


\section*{\refname}
\printbibliography[heading=none]
\end{refsection}

